\begin{helper}[Part (i) of the comparison scheme: the Morrey bound on the middle domain]
Let $n\ge2$, $p>1$, $0<\tau<n$, $c_0>0$, $W\ge0$, $d>0$ and $R_1>0$, and let
$\sigma_2\colon\R\to\R$ be nonnegative, continuous, increasing with
$\sigma_2(0)=0$.  Then there is $N_1\ge0$, depending only on these quantities,
with the following property.

Let $0\le\varepsilon\le1$, let $\Omega,\Omega''\subseteq\R^n$ be open with
$\Omega''$ bounded and $\overline{\Omega''}\subseteq\Omega$, and with the
separation
\[
d\le|x-y|\qquad\text{for }x\in\Omega'',\ y\notin\Omega .
\]
Let $u\in W^{1,p}(\Omega)$ with weak gradient $g=\nabla u$
(\noderef{MemW1pOn}) satisfy
\[
\int_\Omega|u|^p\,dx+\int_\Omega|\nabla u|^p\,dx\le W^p .
\]
Assume finally that for every $x_0\in\Omega''$ and every $0<R\le R_1$ with
$\overline B_R(x_0)\subseteq\Omega$ there is a comparison map
$v\in W^{1,p}\bigl(B_R(x_0)\bigr)$ with weak gradient $g_v$
(\noderef{MemW1pOn}) such that
\begin{align}
\int_{B_\rho(x_0)}(\varepsilon+|g_v|^2)^{\frac p2}dx
&\le c_0\Bigl(\frac\rho R\Bigr)^{n}\int_{B_R(x_0)}(\varepsilon+|g_v|^2)^{\frac p2}dx
\qquad(0<\rho\le R),\tag{4.11}\\
\int_{B_R(x_0)}|\nabla u-g_v|^p\,dx
&\le\sigma_2(R)\int_{B_R(x_0)}(1+|\nabla u|^2)^{\frac p2}dx .\tag{$\star$}
\end{align}
Then
\[
\|\,|\nabla u|\,\|_{L^{p,n-\tau}(\Omega'')}\le N_1
\qquad(\noderef{MorreyBoundOn}).
\]

\noindent\textbf{Why the bound is stated on $\Omega''$ and not on $\Omega'$.}
The thesis' Theorem 3(i) asserts the Morrey bound on the inner domain
$\Omega'$, but its proof produces it on the middle domain $\Omega''$ on which
the comparison bundle lives, and the passage from (i) to (ii) --- Lemma~4.2 at
a ball $B_{d/4}(x_1)$ centred at a point of $\Omega'$ --- consumes the
$\Omega''$ form, not the $\Omega'$ form, because that ball need not lie in
$\Omega'$.  Stating the node on $\Omega''$ therefore records exactly what the
argument gives and what its two consumers need; the $\Omega'$ form follows in
one line by restriction, since $\Omega'\subseteq\Omega''$.

\noindent\textbf{On the $W$ datum.}  The hypothesis bounds
$\int_\Omega|u|^p+\int_\Omega|\nabla u|^p$, i.e.\ the genuine Sobolev norm, and
not $\int_\Omega(1+|\nabla u|^2)^{p/2}dx$.  The latter integrand is pointwise at
least $1$, so a bound on it would force $W\ge|\Omega|$ and would be
unsatisfiable when $|\Omega|=\infty$; since the theorems of \S4 and \S5 that
build this bundle are stated for an arbitrary open $\Omega$, that would make
them vacuous on instances such as $\Omega=\R^n$.  The weighted integral is
still available on each ball, because $(1+t^2)^{p/2}\le2^{p/2}(1+t^p)$ for
$t\ge0$ gives
\[
\int_{B_R}(1+|\nabla u|^2)^{\frac p2}dx
\;\le\;2^{\frac p2}\Bigl(|B_R|+\int_{B_R}|\nabla u|^p\,dx\Bigr)
\;\le\;2^{\frac p2}\bigl(\omega_nR^{\,n}+W^p\bigr),
\]
whose right-hand side involves only $n$, $p$, $R$ and $W$.

\noindent\textbf{On the membership $v\in W^{1,p}(B_R(x_0))$.}  The comparison
map of \S4 and \S5 is a $W^{1,p}$ solution of a frozen Dirichlet problem, so
this is what every producer of the bundle has.  It is recorded explicitly
because (4.11) of the thesis has the \emph{same} integrand on both sides and
$(\star)$ bounds an integral of $|\nabla u-g_v|^p$: without it a field whose
$(\varepsilon+|g_v|^2)^{p/2}$ fails to be integrable on the ball would satisfy
both displays with every integral equal to the junk value, and the hypothesis
would constrain nothing.
\end{helper}
\begin{proof}
Throughout, $B_t=B_t(x_0)$ and $\overline B_t=\overline B_t(x_0)$ for the ball
currently in play, $|E|$ is Lebesgue measure, $\omega_n=|B_1(0)|\in(0,\infty)$
(the Preamble item on the volume of a ball), and all non-integer powers are
real powers of nonnegative numbers.  Put
\[
\kappa=\max\bigl\{1,2^{\frac p2-1}\bigr\},
\]
so that $(X+Y)^{p/2}\le\kappa\bigl(X^{p/2}+Y^{p/2}\bigr)$ for all $X,Y\ge0$: for
$p\le2$ the exponent $p/2$ lies in $(0,1]$ and $\lambda\mapsto\lambda^{p/2}$ is
subadditive on $[0,\infty)$, while for $p\ge2$ it is the convexity of
$\lambda\mapsto\lambda^{p/2}$ in the form
$\bigl(\frac{X+Y}2\bigr)^{p/2}\le\frac12\bigl(X^{p/2}+Y^{p/2}\bigr)$.


\emph{Step 0: two elementary comparisons.}  For $a,b\in\R^n$ and $\varepsilon\ge0$,
\[
\varepsilon+|a|^2\le2\bigl(\varepsilon+|b|^2\bigr)+2|a-b|^2 ,
\]
because $|a|^2\le2|b|^2+2|a-b|^2$ (from $|a|\le|b|+|a-b|$ and
$(s+u)^2\le2s^2+2u^2$) and $\varepsilon\le2\varepsilon$.  Raising to the power
$p/2$ and using $\kappa$,
\[
\bigl(\varepsilon+|a|^2\bigr)^{\frac p2}
\le\kappa\,2^{\frac p2}\Bigl[\bigl(\varepsilon+|b|^2\bigr)^{\frac p2}+|a-b|^{p}\Bigr]
=:c_1\Bigl[\bigl(\varepsilon+|b|^2\bigr)^{\frac p2}+|a-b|^{p}\Bigr],
\tag{E1}
\]
with $c_1=c_1(p)=\kappa2^{p/2}\ge1$, where $\bigl(|a-b|^2\bigr)^{p/2}=|a-b|^p$.
Interchanging $a$ and $b$ gives the same inequality with the two roles
exchanged, which is the form used twice below.  Second, for $0\le\varepsilon\le1$,
\[
\bigl(1+|a|^2\bigr)^{\frac p2}
\le\bigl(\bigl(\varepsilon+|a|^2\bigr)+1\bigr)^{\frac p2}
\le\kappa\Bigl[\bigl(\varepsilon+|a|^2\bigr)^{\frac p2}+1\Bigr],
\tag{E2}
\]
the first step because $1+|a|^2\le1+\varepsilon+|a|^2$.

\emph{Step 1: the constants, chosen before all the data.}  Put
\[
c_2=c_0c_1^2,\qquad
c_3=c_0c_1^2+c_1,\qquad
c_3'=\kappa\,c_3,\qquad
A=\max\{c_2,c_3'\} ,
\]
all positive and determined by $n$, $p$ and $c_0$ alone.  Apply
\noderef{CampanatoIterationLemma} with this $A$, with $\alpha=n$ and with
$\beta=n-\tau$: its hypotheses $0<A$, $0<\alpha$, $0<\beta$ and $\beta<\alpha$
hold, the last two because $0<\tau<n$.  It supplies
$\delta_0=\delta_0(A,n,n-\tau)>0$ and $C=C(A,n,n-\tau)>0$.  Since $\sigma_2$ is
continuous with $\sigma_2(0)=0$, there is $R_2>0$ with
\[
\sigma_2(R_2)<\delta_0 ,
\]
and we fix one such $R_2$.  Finally put
\[
c_4=c_3'\bigl(1+\sigma_2(R_2)\bigr)\,\omega_n\bigl(1+R_2^{\,\tau}\bigr)>0,
\qquad
c_5=C>0 .
\]
All of $c_1,\kappa,c_2,c_3,c_3',A,\delta_0,C,R_2,c_4,c_5$ depend only on
$n,p,\tau,c_0$ and $\sigma_2$.  The radius and the constant of the conclusion
are fixed here, from these and from the remaining parameters $W$, $d$ and
$R_1$:
\begin{enumerate}
\item[(1e)] put
  \[
  R_\ast=\min\Bigl\{R_1,\ R_2,\ \frac d{16}\Bigr\}>0,
  \qquad
  A_1=c_5R_\ast^{-(n-\tau)}\Bigl\{2^{\frac p2}\bigl(\omega_nR_\ast^{\,n}+W^p\bigr)
      +c_4R_\ast^{\,n-\tau}\Bigr\},
  \]
  \[
  N_1=\Bigl[\max\bigl\{A_1,\ W^pR_\ast^{-(n-\tau)}\bigr\}\Bigr]^{1/p}\ \ge0 ;
  \]
\end{enumerate}
Every constant named so far is built from $n,p,\tau,c_0,W,d,R_1$ and $\sigma_2$
alone --- in particular not from $\varepsilon,\Omega,\Omega'',u,g$ --- which is
the quantifier order the statement prescribes.  Now fix those data satisfying
all the hypotheses.

\emph{Step 2: the setting for the rest of the proof.}  Let $x_0\in\Omega''$ and
let $0<\rho\le R\le\min\{R_1,R_2\}$ with $\overline B_R\subseteq\Omega$.  Define
\[
\varphi\colon\R\to\R,\qquad
\varphi(t)=\int_{B_{\min\{t,R\}}}\bigl(\varepsilon+|g|^2\bigr)^{\frac p2}dx
\quad(t>0),\qquad \varphi(t)=0\quad(t\le0),
\]
the radius being capped at $R$ so that $\varphi$ is defined, finite and
nondecreasing on all of $\R$, as the iteration lemma requires.  Indeed:
$\varphi\ge0$ because the integrand is a real power of a nonnegative number;
$\varphi$ is nondecreasing because $s\le t$ gives
$B_{\min\{s,R\}}\subseteq B_{\min\{t,R\}}$ and the integrand is nonnegative; and
$\varphi$ is finite because for $t\le R$, using $\varepsilon\le1$ and
$(1+|h|^2)^{p/2}\le2^{p/2}(1+|h|^{p})$,
\[
\varphi(t)\le\int_{B_R}\bigl(1+|g|^2\bigr)^{\frac p2}dx
\le2^{\frac p2}\Bigl(|B_R|+\int_{B_R}|g|^{p}dx\Bigr)
\le2^{\frac p2}\bigl(\omega_nR^{\,n}+W^{p}\bigr)<\infty ,
\tag{F}
\]
where $\int_{B_R}|g|^p\le\int_\Omega|g|^p\le W^p$, the last step because
$\int_\Omega|u|^p+\int_\Omega|g|^p\le W^p$ and the first of those two integrals
is nonnegative, and
monotonicity of the integral of a nonnegative function over subsets (Preamble),
legitimate because $B_R\subseteq\overline B_R\subseteq\Omega$.  Note for later that $\varphi(t)=\int_{B_t}(\varepsilon+|g|^2)^{p/2}dx$
whenever $0<t\le R$.

\emph{Step 3: the passage to (4.12), written out.}  Let $0<\rho'\le R'\le R$.
Then $R'\le R_1$ and $\overline B_{R'}\subseteq\overline B_R\subseteq\Omega$, so
the hypothesis of the lemma applies at the radius $R'$ and supplies a field
$g_v$ satisfying (4.11) and $(\star)$ on $B_{R'}$.  Write $|\nabla w|=|g-g_v|$
for the remainder, as in the thesis.  All the integrals written below are finite:
each integrand is nonnegative and dominated on $B_{R'}\subseteq\Omega$ either by
$(1+|g|^2)^{p/2}$, which is integrable there by $(F)$, or --- in the case of
$(\varepsilon+|g_v|^2)^{p/2}$ and $|\nabla w|^p$ --- by the right-hand sides of
(4.11) and $(\star)$, which are finite by $(F)$ and the hypothesis $(\star)$
itself.

First, by (E1) with $a=g$, $b=g_v$, integrated over $B_{\rho'}$ and then
enlarging the domain of the second term to $B_{R'}$ (the integrand being
nonnegative and $\rho'\le R'$),
\[
\int_{B_{\rho'}}\bigl(\varepsilon+|g|^2\bigr)^{\frac p2}dx
\ \le\ c_1\int_{B_{\rho'}}\bigl(\varepsilon+|g_v|^2\bigr)^{\frac p2}dx
+c_1\int_{B_{R'}}|\nabla w|^{p}dx .
\tag{3a}
\]
Second, (4.11) at the pair $(\rho',R')$ bounds the first term:
\[
\int_{B_{\rho'}}\bigl(\varepsilon+|g_v|^2\bigr)^{\frac p2}dx
\ \le\ c_0\Bigl(\frac{\rho'}{R'}\Bigr)^{n}
\int_{B_{R'}}\bigl(\varepsilon+|g_v|^2\bigr)^{\frac p2}dx .
\tag{3b}
\]
Third --- and this is the step where the comparison is returned to $u$ --- (E1)
with the roles interchanged, i.e.\ with $a=g_v$ and $b=g$, integrated over
$B_{R'}$, gives
\[
\int_{B_{R'}}\bigl(\varepsilon+|g_v|^2\bigr)^{\frac p2}dx
\ \le\ c_1\int_{B_{R'}}\bigl(\varepsilon+|g|^2\bigr)^{\frac p2}dx
+c_1\int_{B_{R'}}|\nabla w|^{p}dx .
\tag{3c}
\]
Substituting (3c) into (3b) and then into (3a), and using
$(\rho'/R')^{n}\le1$ on the term carrying $|\nabla w|^p$,
\[
\int_{B_{\rho'}}\bigl(\varepsilon+|g|^2\bigr)^{\frac p2}dx
\ \le\ c_0c_1^2\Bigl(\frac{\rho'}{R'}\Bigr)^{n}
\int_{B_{R'}}\bigl(\varepsilon+|g|^2\bigr)^{\frac p2}dx
+\bigl(c_0c_1^2+c_1\bigr)\int_{B_{R'}}|\nabla w|^{p}dx .
\tag{3d}
\]
Now $(\star)$ is applied --- and this is the second of the two applications the
thesis means by ``using a second time Lemma 4.1'': the first is the one that
disposes of the remainder produced by (3a) on the small ball, the second is the
one that disposes of the remainder produced by (3c) when the $v$-integral over
$B_{R'}$ is returned to $u$.  Both remainders are the same quantity
$\int_{B_{R'}}|\nabla w|^p$, so the two applications may be performed at once on
the combined coefficient $c_3=c_0c_1^2+c_1$:
\[
\bigl(c_0c_1^2+c_1\bigr)\int_{B_{R'}}|\nabla w|^{p}dx
\ \le\ c_3\,\sigma_2(R')\int_{B_{R'}}\bigl(1+|g|^2\bigr)^{\frac p2}dx .
\]
With (3d) this is the analogue of (4.12) of the thesis for the pair $(\rho',R')$:
\[
\int_{B_{\rho'}}\bigl(\varepsilon+|g|^2\bigr)^{\frac p2}dx
\ \le\ c_2\Bigl(\frac{\rho'}{R'}\Bigr)^{n}
\int_{B_{R'}}\bigl(\varepsilon+|g|^2\bigr)^{\frac p2}dx
+c_3\,\sigma_2(R')\int_{B_{R'}}\bigl(1+|g|^2\bigr)^{\frac p2}dx .
\tag{4.12$'$}
\]

\emph{Step 4: the passage to (4.13).}  By (E2) integrated over $B_{R'}$,
\[
\int_{B_{R'}}\bigl(1+|g|^2\bigr)^{\frac p2}dx
\ \le\ \kappa\int_{B_{R'}}\bigl(\varepsilon+|g|^2\bigr)^{\frac p2}dx+\kappa\,|B_{R'}| ,
\]
and $|B_{R'}|=\omega_nR'^{\,n}=\omega_nR'^{\,\tau}R'^{\,n-\tau}
\le\omega_nR_2^{\,\tau}R'^{\,n-\tau}$ because $0<R'\le R\le R_2$ and $\tau>0$.
Also $\sigma_2(R')\le\sigma_2(R_2)$ because $\sigma_2$ is increasing and
$R'\le R_2$.  Substituting into (4.12$'$) and recalling $\varphi(\rho')$,
$\varphi(R')$ from Step 2 --- legitimate because $\rho'\le R'\le R$, so the cap
is inactive ---
\[
\varphi(\rho')
\ \le\ \Bigl[c_2\Bigl(\frac{\rho'}{R'}\Bigr)^{n}+c_3'\,\sigma_2(R')\Bigr]\varphi(R')
+c_4R'^{\,n-\tau}
\tag{4.13$'$}
\]
for all $0<\rho'\le R'\le R$, where $c_3'=\kappa c_3$ and where the last term is
bounded as
$c_3\kappa\,\sigma_2(R')\,|B_{R'}|
\le c_3'\sigma_2(R_2)\omega_nR_2^{\,\tau}R'^{\,n-\tau}\le c_4R'^{\,n-\tau}$,
by the definition of $c_4$ in Step 1.

\emph{Step 5: the iteration lemma.}  Apply \noderef{CampanatoIterationLemma} ---
the same instance as in Step 1, so the $\delta_0$ and $C$ it supplies are those
fixed there --- with
\[
\delta=\sigma_2(R),\qquad B=c_4,\qquad R_0=R,\qquad \varphi \text{ as in Step 2}.
\]
Its hypotheses hold: $0\le\delta$ because $\sigma_2\ge0$; $\delta<\delta_0$
because $\sigma_2$ is increasing, $R\le R_2$ and $\sigma_2(R_2)<\delta_0$;
$0<B=c_4$ and $0<R_0=R$; $\varphi$ is nonnegative and nondecreasing on $\R$ by
Step 2; and the recursive inequality
\[
\varphi(\rho')\le A\Bigl[\Bigl(\frac{\rho'}{R'}\Bigr)^{n}+\delta\Bigr]\varphi(R')
+BR'^{\,n-\tau}
\qquad(0<\rho'\le R'\le R_0=R)
\]
follows from (4.13$'$), because $A\ge c_2$, $A\ge c_3'$,
$\sigma_2(R')\le\sigma_2(R)=\delta$ and $\varphi(R')\ge0$.  The conclusion of the
lemma, at the pair $(\rho,R)$ with $0<\rho\le R=R_0$, is
\[
\varphi(\rho)\ \le\ C\Bigl(\frac\rho R\Bigr)^{n-\tau}
\bigl[\varphi(R)+c_4R^{\,n-\tau}\bigr],
\]
which, with $c_5=C$ and with $\varphi(\rho)$, $\varphi(R)$ written back as the
integrals they are (Step 2, the cap being inactive at $\rho$ and at $R$), is
(4.14) of the thesis:
\[
\int_{B_\rho(x_0)}\bigl(\varepsilon+|\nabla u|^2\bigr)^{\frac p2}dx
\le c_5\Bigl(\frac\rho R\Bigr)^{n-\tau}
\Bigl\{\int_{B_R(x_0)}\bigl(\varepsilon+|\nabla u|^2\bigr)^{\frac p2}dx
+c_4R^{\,n-\tau}\Bigr\}.
\]
As $x_0$, $\rho$ and $R$ were arbitrary subject to the stated constraints, we
have proved
\[
\int_{B_\rho(x_0)}\bigl(\varepsilon+|\nabla u|^2\bigr)^{\frac p2}dx
\le c_5\Bigl(\frac\rho R\Bigr)^{n-\tau}
\Bigl\{\int_{B_R(x_0)}\bigl(\varepsilon+|\nabla u|^2\bigr)^{\frac p2}dx
+c_4R^{\,n-\tau}\Bigr\}
\tag{4.14}
\]
for every $x_0\in\Omega''$ and all $0<\rho\le R\le\min\{R_1,R_2\}$ with
$\overline B_R(x_0)\subseteq\Omega$.  Steps 6--9 below derive the three
conclusions from it.  (The thesis reaches the same display through
(4.13) of the thesis, whose printed right-hand side carries the misprint
$c_3(\tau_2(R))$ for $c_3\sigma_2(R)$; the version proved here is the corrected
one, and the additive term $c_4R^{n-\tau}$ is exactly the thesis'.)
\emph{Step 6: part (i).}  We first prove
$\|\,|g|\,\|_{L^{p,n-\tau}(\Omega'')}\le N_1$
(\noderef{MorreyBoundOn}).  The $L^p$ clause holds: $|g|\in L^p(\Omega)$ by
\noderef{MemW1pOn}, hence $|g|\in L^p(\Omega'')$ by the Preamble
transport fact for $L^q$ membership under restriction to a subset.  Now let
$x_1\in\Omega''$ and $0<\rho<\operatorname{diam}\Omega''$.
Pointwise $|g|^p=(|g|^2)^{p/2}\le(\varepsilon+|g|^2)^{p/2}$, because
$t\mapsto t^{p/2}$ is nondecreasing on $[0,\infty)$ and $\varepsilon\ge0$, and
$\Omega''\cap B_\rho(x_1)\subseteq B_\rho(x_1)$; we distinguish two
cases.

If $\rho\le R_\ast$: by $(4.14)$ with $R=R_\ast$ --- legitimate because
$R_\ast\le d/16<d$ and $x_1\in\Omega''$, so that every $z$ with
$|z-x_1|\le R_\ast$ lies in $\Omega$ by the hypothesis that
$d\le|x-y|$ for every $x\in\Omega''$ and every $y\notin\Omega$ ---
and by $(F)$ of Step 2 together with
$|B_{R_\ast}|=\omega_nR_\ast^{\,n}$ and $\int_\Omega|g|^p\le W^p$,
\[
\int_{\Omega''\cap B_\rho(x_1)}|g|^p\,dx
\le\int_{B_\rho(x_1)}(\varepsilon+|g|^2)^{\frac p2}dx
\le c_5\frac{\rho^{\,n-\tau}}{R_\ast^{\,n-\tau}}
\Bigl\{2^{\frac p2}\bigl(\omega_nR_\ast^{\,n}+W^p\bigr)+c_4R_\ast^{\,n-\tau}\Bigr\}
=A_1\rho^{\,n-\tau},
\]
with $A_1$ as in (1e); here
$\int_{B_{R_\ast}}(\varepsilon+|g|^2)^{p/2}\le2^{p/2}(|B_{R_\ast}|+\int_{B_{R_\ast}}|g|^p)
\le2^{p/2}(\omega_nR_\ast^n+W^p)$ by the pointwise inequality used in $(F)$ of Step 2.

If $\rho>R_\ast$: since $n-\tau>0$ we have
$\rho^{-(n-\tau)}<R_\ast^{-(n-\tau)}$, so
\[
\int_{\Omega''\cap B_\rho(x_1)}|g|^p\,dx\le\int_\Omega|g|^p\,dx\le W^p
=W^p\rho^{-(n-\tau)}\rho^{\,n-\tau}\le W^pR_\ast^{-(n-\tau)}\rho^{\,n-\tau}.
\]
In both cases the left-hand side is at most $N_1^p\rho^{\,n-\tau}$ with $N_1$ as
in (1e).  So $\|\,|g|\,\|_{L^{p,n-\tau}(\Omega'')}\le N_1$.

\end{proof}
